% Adapted from ../zju-tech-report/zju-tech-report.tex (technical report mode).
% Compile directly with XeLaTeX. Build instructions are in README.md.
\documentclass[UTF8,11pt,a4paper,twoside,fontset=none,heading=false]{ctexart}

\usepackage[left=22mm,right=22mm,top=15mm,bottom=21mm,includehead,headheight=14pt,headsep=6mm,footskip=12mm]{geometry}
\setmainfont{TeX Gyre Termes}
\setsansfont{Latin Modern Sans}
\setmonofont{Latin Modern Mono}
\setCJKmainfont{SimSun}[BoldFont=SimHei]
\setCJKsansfont{SimHei}
\usepackage{amsmath,array,booktabs,calc,caption,enumitem,fancyhdr,graphicx,pgfplots,tabularx,titlesec,xcolor}
\usepackage{hyperref}
\pgfplotsset{compat=1.18}

% Qiushi Blue follows the repository template.
\definecolor{QiushiBlue}{RGB}{0,63,136}
\colorlet{ReportBlue}{QiushiBlue}
\definecolor{TextGray}{HTML}{555B66}
\newcommand{\CourseAssignmentTitle}{2026秋季班：写作课程作业}
\hypersetup{
  colorlinks=true,
  linkcolor=ReportBlue,
  urlcolor=ReportBlue,
  pdftitle={\CourseAssignmentTitle},
  pdfauthor={蒋佳豪（22671064）}
}
\newcommand{\ReportDate}{2026年12月8日}

\titleformat{\section}
  {\rmfamily\bfseries\fontsize{14}{18}\selectfont\color{ReportBlue}}
  {\thesection}{0.65em}{}
\titleformat{\subsection}
  {\rmfamily\bfseries\fontsize{12}{16}\selectfont\color{ReportBlue}}
  {\thesubsection}{0.65em}{}
\titlespacing*{\section}{0pt}{1.8ex plus 0.5ex minus 0.2ex}{0.7ex}
\titlespacing*{\subsection}{0pt}{1.3ex plus 0.4ex minus 0.2ex}{0.5ex}
\captionsetup{font=small,labelfont=bf,skip=5pt,hypcap=false}
\setlist{itemsep=5pt,topsep=4pt,leftmargin=1.6em}
\providecommand{\tightlist}{}

\linespread{1.12}
\setlength{\parindent}{0pt}
\setlength{\parskip}{5pt}
\setlength{\emergencystretch}{2em}
\clubpenalty=10000
\widowpenalty=10000
\setcounter{secnumdepth}{2}

\newcommand{\makecoursetitle}{%
  {\centering\fontsize{18}{24}\selectfont
    \CourseAssignmentTitle\par}
  \vspace{12pt}
  {\fontsize{12}{17}\selectfont
    课程名称：\textbf{科学研究与写作指导，开\ 课\ 号：1147}\par
    \textbf{开课院系：国际联合学院（海宁国际校区）ZJUI，学期：2026-2027学年秋冬学期}\par
    \textbf{开课教师：李尔平教授，\mbox{liep@zju.edu.cn}}\par}
  \vspace{12pt}
  {\centering 蒋佳豪（22671064）\quad 提交日期：\ReportDate\par}
  \vspace{11pt}
  \textbf{报告提要}\quad
  本报告分析 Mon-Williams 等人的 ELLMER 论文，逐项检查摘要、引言和正文如何提出问题、说明方法并报告结果。真实机械臂演示支持该方案在实验场景中的可行性。规划文本的对比结果进一步支持检索增强生成的作用。论文可以通过补充整项任务的统计、说明基线设置和限定结论范围，使主张与证据更容易对应。\par
  \vspace{8pt}
  \thispagestyle{plain}
}
\pagestyle{fancy}
\fancyhf{}
\fancyhead[RO,RE]{\footnotesize\rmfamily\color{TextGray}\CourseAssignmentTitle}
\fancyfoot[C]{\footnotesize\rmfamily\color{TextGray}\thepage}
\renewcommand{\headrulewidth}{0pt}
\renewcommand{\footrulewidth}{0pt}
\fancypagestyle{plain}{%
  \fancyhf{}
  \fancyfoot[C]{\footnotesize\rmfamily\color{TextGray}\thepage}
  \renewcommand{\headrulewidth}{0pt}
  \renewcommand{\footrulewidth}{0pt}
}

\begin{document}
\makecoursetitle
\section{参考论文信息}

\textbf{题目}　Embodied large language models enable robots to complete complex tasks in unpredictable environments

\textbf{作者}　Ruaridh Mon-Williams（1、2、3）、Gen Li（1）、Ran Long（1）、Wenqian Du（1、4）、Christopher G. Lucas（1）

\textbf{作者单位}　（1）英国爱丁堡大学，（2）美国麻省理工学院，（3）美国普林斯顿大学，（4）英国艾伦·图灵研究所

\textbf{期刊与发表时间}　\emph{Nature Machine Intelligence}，第 7 卷，592--601 页。2025 年 3 月 19 日在线发表，收入 2025 年 4 月刊。\href{https://doi.org/10.1038/s42256-025-01005-x}{DOI 10.1038/s42256-025-01005-x}

\section{摘要的四个要点}

以下四项按原文摘要（第 592 页）的内容判断。

\subsection{研究问题是什么}

\textbf{已说明。} 机器人需要在环境发生变化时完成连续操作。摘要将难点引向语言理解与感觉运动控制的结合，为后文同时引入大语言模型和传感器反馈提供了理由。摘要未举例说明哪些环境变化由系统处理，读者需到引言才知道杯子被移动、抽屉开启方式未知等具体情形。

\subsection{过去做了什么}

\textbf{未具体说明。} 摘要提到人类的感觉运动能力和受生物启发的机器智能。这些内容交代了研究动机。已有系统如何把语言指令转成动作，以及它们在哪些接触任务中仍遇到困难，留到了引言。

\subsection{本文有什么新意}

\textbf{已说明。} ELLMER 用 GPT-4 理解指令，再用检索增强生成（RAG）从知识库选取相关的动作示例。生成的计划调用视觉和力反馈来调整执行。摘要由此交代了系统贡献，即把语言规划、示例检索和感知控制接入同一套机器人系统。

\subsection{主要发现是什么}

\textbf{已报告演示结果。} 机械臂完成了制作咖啡和装饰盘子。摘要点出开抽屉、倒水等子任务需要不同反馈，但没有报告重复次数、整项任务完成率或对照结果。可在结果句补入一个有明确条件的指标，例如第 595 页的白杯遮挡测试，并说明它衡量的是识别性能。

四项中，摘要缺少对既有研究的具体交代。主要发现已有定性描述，量化证据可进一步补充。

\section{引言的写作评析}

\subsection{引言有没有说清研究问题}

\textbf{基本说清。} 开篇以深蓝能否自行移动棋子引出具身认知，随后用厨房场景说明工程问题。机器人收到高层指令后，要找杯子、判断抽屉如何打开、计量倒水，还需在杯子被移动时继续执行。这个例子让读者看到研究涉及任务分解、物体定位和接触控制。引言末尾给出了机器人能否为人制作咖啡的可检验问题。开头对“类人智能”的讨论范围更广，可在提出实验时明确本研究考察的是哪些可观察的操作能力（第 592--593 页）。

\subsection{已有哪些方法解决这个问题}

\textbf{交代了主要路线。} 引言提到强化学习可通过交互获得动作策略，模仿学习可从示范中学习操作，也介绍了用语言模型规划动作、用视觉语言模型迁移知识和用分层策略生成轨迹的路线。作者由此指出，长流程任务还需要处理执行中的环境扰动和接触力变化，并提出把可检索的动作示例与视觉和力反馈结合。这里可以更清楚地说明哪些已有系统已能处理扰动，以及本文新增了哪类反馈和任务条件（第 593 页）。

\subsection{引文是否说明了各项工作的具体贡献}

\textbf{部分符合。} 原文逐项介绍 VoxPoser、RT-2 和层级扩散策略，其余若干判断用成组编号支撑。以下沿用原论文编号，说明关键引文与本文论证的关系。

\begin{itemize}
\tightlist
\item
  \textbf{{[}15{]} VoxPoser（Huang 等）}用语言模型推断操作目标与约束，再通过三维价值图生成机器人轨迹。该系统已有闭环操作能力。与 ELLMER 对比时，应具体讨论视觉与力反馈如何进入各自的执行过程。\href{https://proceedings.mlr.press/v229/huang23b.html}{会议论文}
\item
  \textbf{{[}27{]} 零样本语言规划（Huang 等）}从预训练语言模型中提取任务步骤，并将这些步骤映射到环境中可执行的动作。该工作为本文使用语言模型分解任务提供先例。\href{https://proceedings.mlr.press/v162/huang22a.html}{会议论文}
\item
  \textbf{{[}29{]} RT-2（Zitkovich 等）}用网络视觉语言数据和机器人轨迹共同微调模型，将知识迁移到机器人控制。它为陌生物体和新指令的适应能力提供证据，代表一种端到端路线。\href{https://proceedings.mlr.press/v229/zitkovich23a.html}{会议论文}
\item
  \textbf{{[}30{]} 层级扩散策略（Ma 等）}由高层预测下一个末端执行器位姿，再由低层扩散策略生成符合运动学约束的轨迹。它提供了分层控制的参照。描述这项前作时，应明确高层输出的是目标位姿。\href{https://openaccess.thecvf.com/content/CVPR2024/html/Ma_Hierarchical_Diffusion_Policy_for_Kinematics-Aware_Multi-Task_Robotic_Manipulation_CVPR_2024_paper.html}{会议论文}
\item
  \textbf{{[}32{]} 检索增强生成（Lewis 等）}在语言生成时从外部语料检索相关内容。ELLMER 借用了这种机制，将检索对象改为机器人动作的代码示例。\href{https://proceedings.nips.cc/paper/2020/hash/6b493230205f780e1bc26945df7481e5-Abstract.html}{会议论文}
\end{itemize}

第 593 页用 {[}16{]} 至 {[}23{]} 支撑模仿学习的适应性限制，又以 {[}36{]} 至 {[}53{]} 指向补充背景，正文未逐篇说明这些引文的作用。比如，{[}16{]} 研究机器人倒饮料，{[}19{]}（CLIPort）处理语言指令与空间操作位置的对应。用一个句子分别交代任务、成果和实验边界，能让读者看出各文献支持的是哪一项判断。

\subsection{引言是否按章节预告全文结构}

\textbf{缺少章节导览。} 引言末段提出 ELLMER 和咖啡制作实验，并指向图 1。图 1 提供系统概览。若补写篇章安排，可以说明 Results 报告任务演示和模块评估，Discussion 解释适用条件，Methods 交代硬件、检索与控制实现。这样的导览也能帮助读者在先见到结果、后读到方法的期刊结构中找到评价细节。

\section{正文的四个问题}

\subsection{正文要解决什么问题}

作者要让一台机械臂把开放式的口头请求分解成有先后依赖的动作，并根据现场观测调整这些动作。以制作热饮为例，杯子能否直接找到会决定是否插入开抽屉的动作。倒水时杯位会影响末端轨迹，液体重量又决定何时停止。图 1 将这一问题落实为图像与语言输入、动作代码生成、机械臂执行和传感器反馈（第 593、597--598 页）。

\subsection{作者为什么选择这些方法}

\textbf{用语言模型分解任务。} GPT-4 根据指令和初始图像选择任务并生成代码。作者在方法节注明，可执行控制代码的思路继承自 Code as Policies {[}65{]}，该工作已用语言模型生成能够接收感知反馈的机器人程序。\href{https://ieeexplore.ieee.org/abstract/document/10160591/}{会议论文}

\textbf{用 RAG 检索动作示例。} 系统按请求检索开抽屉、舀取、倒水等已有函数，使模型可以组合并调整经过验证的动作。作者选用自定义 GPT 的内置检索功能，理由是实现方便，而且在其任务设置中能稳定提供相关示例。知识库扩大后的检索质量与计算代价仍需单独测量（第 594、598 页）。

\textbf{用视觉定位、用力反馈控制接触。} RGB-D 视觉估计物体位姿，帮助机械臂重新定位被移动的杯子。力和力矩传感器提供接触信息，可用于判断放杯、估计倒水量和控制画笔压力。ROS 控制器接收生成的代码并执行反馈控制，感知更新速度因模块而异（图 1，第 595、598 页）。

\subsection{作者做了哪些假设}

\textbf{感知与动作先验。} 讨论节列出两项前提，即视觉模块能正确识别物体，并且系统掌握器具的可供性，亦即器具可以怎样被抓取或操作。实验给定了水壶、勺子和门把手的相关信息。知识库也预先包含主要动作示例（第 594、597 页）。

\textbf{低层控制与实验范围。} 方法节假设低层控制负责避障，并在动作函数中设置速度、力和工作空间限制。实验使用一台机械臂和一组预设的厨房物件。这些条件决定了文中结论首先适用于所测试的环境和扰动（第 597--599 页）。

\textbf{倒水的准静态近似。} 作者假设倒水足够缓慢，传感器测得的竖直力变化主要来自液体质量变化。加速度增大后，惯性和质心变化会进入测量值。第 595 页报告的高速条件下误差增加，与该近似的适用范围相符（第 595、598 页）。

\clearpage
\subsection{作者如何验证结果}
\label{sec:evaluation}

作者在三个层面检查 ELLMER：先演示整套机械臂流程，再测量感知和力反馈的表现，最后用对照实验评价语言模型生成的计划。三组结果回答的问题不同，表~\ref{tab:evidence}列出了它们各自能支持的结论。

\textbf{整项任务演示。} 原文图 2、图 6 和补充视频 1 记录了机械臂从接收指令到制作咖啡和装饰盘子的过程。开抽屉、取杯、舀取、倒水及绘画等步骤需要不同的感知或接触信息。这项演示说明 ELLMER 可以在所搭建的场景中串联这些技能。原文没有汇报整项任务的重复次数、完成率或失败分布，因而还无法估计重复执行的可靠性（第 593--594、597 页）。

\textbf{模块测量。} 图~\ref{fig:results}(a) 汇总两档白杯遮挡条件。遮挡在 20\% 到 30\% 时，原文报告平均识别率约为 90\%。遮挡升至 80\% 到 90\% 时，识别率约为 20\%。原文图 4 同时给出了不同操作阶段的力、速度和位置曲线。低速倾倒条件下，每 100 克液体的计量误差约为 5.4 克。随着倾倒速度增加，原文报告误差增大，这与方法部分对准静态近似适用范围的讨论相符（第 595、598 页）。

\textbf{生成计划的对照。} 作者用语言模型生成了 80 条覆盖知识库任务范围的请求，比较动态检索和把知识库静态放入上下文两种条件。原文把静态条件称为 VoxPoser 基线。在图~\ref{fig:results}(b) 所列的三个模型中，GPT-4 的忠实度分数增幅最大，由 0.74 升至 0.88。GPT-3.5-turbo 和 Zephyr-7B-beta 分别提高 0.08 和 0.07。这个指标衡量生成回答的忠实度。要判断增益能否转化为更多实物任务成功，还需要固定机械臂、任务与扰动条件后比较整项任务完成率（第 595--596 页）。

\begin{center}
\begin{minipage}{\linewidth}
\centering
\begin{minipage}[t]{0.40\linewidth}
\centering
\textbf{(a)\quad 白杯遮挡测试}\par\smallskip
\begin{tikzpicture}
\begin{axis}[
  scale only axis, width=46mm, height=46mm,
  xmin=0.55, xmax=2.45, ymin=0, ymax=100,
  ybar, bar width=13pt,
  xtick={1,2}, xticklabels={20\%--30\%,80\%--90\%},
  ytick={0,25,50,75,100},
  xlabel={遮挡比例}, ylabel={平均识别率（\%）},
  tick align=outside, tick style={draw=TextGray},
  tick label style={font=\small}, label style={font=\small},
  axis x line*=bottom, axis y line*=left,
  ymajorgrids, grid style={draw=gray!25},
  clip=false
]
\addplot[draw=QiushiBlue, fill=QiushiBlue!75] coordinates {(1,90) (2,20)};
\node[font=\small, anchor=south] at (axis cs:1,90) {约 90\%};
\node[font=\small, anchor=south] at (axis cs:2,20) {约 20\%};
\end{axis}
\end{tikzpicture}
\end{minipage}\hfill
\begin{minipage}[t]{0.57\linewidth}
\centering
\textbf{(b)\quad 计划生成对照}\par\smallskip
\begin{tikzpicture}
\begin{axis}[
  scale only axis, width=54mm, height=46mm,
  xmin=0, xmax=1, ymin=0.5, ymax=3.5,
  xtick={0,0.25,0.50,0.75,1},
  xticklabels={0,0.25,0.50,0.75,1.00},
  ytick={1,2,3}, yticklabels={Zephyr-7B-beta,GPT-3.5-turbo,GPT-4},
  xlabel={回答忠实度},
  tick align=outside, tick style={draw=TextGray},
  tick label style={font=\small}, label style={font=\small},
  axis x line*=bottom, axis y line*=left,
  xmajorgrids, grid style={draw=gray!25},
  legend style={at={(0.5,-0.32)},anchor=north,legend columns=2,
    draw=none,font=\small,column sep=8pt,cells={anchor=west}},
  clip=false
]
\addplot[gray!65, thick, forget plot] coordinates {(0.74,3) (0.88,3)};
\addplot[gray!65, thick, forget plot] coordinates {(0.78,2) (0.86,2)};
\addplot[gray!65, thick, forget plot] coordinates {(0.37,1) (0.44,1)};
\addplot[only marks, mark=o, mark size=2.6pt, thick, black]
  coordinates {(0.74,3) (0.78,2) (0.37,1)};
\addlegendentry{静态上下文}
\addplot[only marks, mark=square*, mark size=2.5pt, thick, QiushiBlue]
  coordinates {(0.88,3) (0.86,2) (0.44,1)};
\addlegendentry{动态检索}
\node[font=\small, anchor=south east] at (axis cs:0.74,3) {0.74};
\node[font=\small, anchor=south west] at (axis cs:0.88,3) {0.88};
\node[font=\small, anchor=south east] at (axis cs:0.78,2) {0.78};
\node[font=\small, anchor=south west] at (axis cs:0.86,2) {0.86};
\node[font=\small, anchor=south east] at (axis cs:0.37,1) {0.37};
\node[font=\small, anchor=south west] at (axis cs:0.44,1) {0.44};
\end{axis}
\end{tikzpicture}
\end{minipage}

\captionof{figure}{白杯识别率与计划生成忠实度}
\label{fig:results}
{\small\raggedright 据 Mon-Williams 等（2025，第 595--596 页）所报数值绘制。(a) 为两组近似识别率。(b) 的连线用于配对同一模型的两种设置。原文相应段落未报告各估计值的不确定性。\par}
\end{minipage}
\end{center}

\begin{center}
\begin{minipage}{\linewidth}
\centering
\captionof{table}{三类证据支持的结论与验证缺口}
\label{tab:evidence}
\small
\renewcommand{\arraystretch}{1.15}
\begin{tabularx}{\linewidth}{
  @{}>{\raggedright\arraybackslash\bfseries}p{0.15\linewidth}
  >{\raggedright\arraybackslash}p{0.36\linewidth}
  >{\raggedright\arraybackslash}X@{}}
\toprule
证据类型 & \textbf{可以得出的结论} & \textbf{还需要的证据} \\
\midrule
实物演示 & 机械臂在所设场景中完成了咖啡制作和盘子装饰。
  & 重复整项任务，报告成功次数、试验总数和失败环节。 \\
\addlinespace[3pt]
模块测量 & 重度遮挡降低白杯识别率。低速倾倒可用力反馈计量。
  & 报告各遮挡组试验数，以及不同速度和容器下的倾倒误差。 \\
\addlinespace[3pt]
规划对照 & 在 80 条生成请求中，动态检索条件的回答忠实度更高。
  & 说明评分与基线实现，比较两种设置的实物任务完成率。 \\
\bottomrule
\end{tabularx}
\end{minipage}
\end{center}

\clearpage
\section{写作上的不足与修改建议}

\subsection*{把关键术语写成可观察条件}
标题和摘要使用“不可预测环境”，引言还讨论“类人智能”。建议在引言用杯子移动、目标遮挡和抽屉开合方式说明实验涵盖的变化，再定义各子任务和整项任务怎样算完成。这样读者可以按同一标准检查结果与结论。

\subsection*{把对照实验的设置交代完整}
正文将静态上下文基线称为 VoxPoser，却未在该段说明复用了哪些原方法组件。建议报告基线与 ELLMER 的共同输入，以及提示词、知识库内容、检索设置和忠实度评分方法。计划忠实度、视觉识别率和实物任务完成率应各自对应具体的评价过程。

若要把三类证据连起来，可在同一机械臂和任务集上运行两种检索设置，并固定提示词、底层控制与扰动条件。每次试验同时记录生成计划的忠实度、实物任务是否完成，以及失败发生的环节。这样的配对设计可以检验忠实度提升是否伴随实物完成率改善。

\subsection*{用试验细节支持量化表述}
原文图 4 的力曲线有助于解释不同操作的接触特征。若同时报告重复次数、测量误差和失败类型，读者便可判断这条曲线是否具有代表性。第 595 页用线速度单位 m/s 表示“pitch velocity”，建议核对速度的物理量、记录位置和单位，并说明它如何影响倒水误差。

\subsection*{将结论收束到数据覆盖的范围}
原文在 Results 和 Discussion 交代了重度遮挡、高速倾倒和中途切换任务的限制。摘要和结论可明确写出已有演示的设备、任务和扰动条件。可扩展性和执行效率则分别需要随知识库规模变化的测试，以及与其他系统在同一任务上的耗时对照。

\section{原论文附件}

期刊发表版全文随附，文件名为 ELLMER\_Nature\_Machine\_Intelligence\_2025.pdf。页码、“原文图”的图号和方括号内的文献编号均沿用该论文。
\end{document}
